% ===========================================================================
\section{Arithmetizing the SHA-256 building blocks}\label{sec:arith}
% ===========================================================================

We collect the lemmas that translate each SHA-256 building block into
either a $\Q[X]$- or an $\F_2[X]$-constraint. The lemmas are stated in
slight variations from the main Zinc+ paper to fit the all-$\Q[X]$
instantiation actually deployed in the implementation: we lift the
$\F_2[X]$ identities back to $\Q[X]$ using extra slack witnesses
whenever a column is not yet bit-poly-typed.

\subsection{Rotations}

We start from the rotation-as-ideal-membership lemma of the main paper
(Lemma~8.7 there), which we reproduce for self-containment.

\begin{lemma}[Rotation as ideal membership]\label{lem:rot}
Let $R$ be a commutative ring and $0 \le r < 32$. For any
$u,v \in R^{<32}[X]$,
$$
v - X^{32-r} \cdot u \;\in\; (X^{32}-1)
\quad \Longleftrightarrow \quad
v \;=\; \ROTR^{r}(u).
$$
\end{lemma}

\subsection{The functions $\Sigma_0$ and $\Sigma_1$}\label{sec:big-sigma}

Define the rotation-encoding polynomials
\begin{equation}\label{eq:rho}
\rho_{0}(X) := X^{30} + X^{19} + X^{10}, \qquad
\rho_{1}(X) := X^{26} + X^{21} + X^{7}, \qquad \text{in } \F_2[X].
\end{equation}

\begin{lemma}[{$\Q[X]$ lift of $\Sigma_0/\Sigma_1$}]\label{lem:big-sigma}
Let $a \in [0\mathbin{..}2^{32}-1]$ with bit-polynomial representative
$\bp{a} \in \bitpoly{32}$, and let $\bp{y} \in \bitpoly{32}$. Then
$\bp{y} = \widehat{\Sigma_0(a)}$ if and only if there exists
$\bp{o} \in \bitpoly{32}$ such that
\begin{equation}\label{eq:sig0-Q}
\bp{a} \cdot \rho_0(X) \;-\; \bp{y} \;-\; 2\,\bp{o} \;\in\; (X^{32}-1)
\qquad \text{in } \Q[X],
\end{equation}
and analogously for $\Sigma_1$ with $\rho_1(X)$.
\end{lemma}

The slack witness $\bp{o}$ absorbs the integer overflow incurred by
lifting the $\F_2[X]$ identity to $\Q[X]$: each coefficient of
$\bp{a}\cdot\rho_0(X) \bmod (X^{32}-1)$ is the sum of three bits, so the
$\Q[X]$-residue $\bp{a}\cdot\rho_0 - \widehat{\Sigma_0(a)}$ has
per-coefficient entries in $\{0,2\}$ and admits a unique bit-poly slack
$\bp{o}$ as in \eqref{eq:sig0-Q}.

\subsection{The functions $\sigma_0$ and $\sigma_1$}\label{sec:small-sigma}

For $\sigma_i$ the right-shift terms are not pure rotations; we model
them with a virtual coordinate-wise map.

\begin{lemma}[Right shift as a virtual column]\label{lem:shr-virtual}
For any commutative ring~$R$ and $r \in \{1,\dots,31\}$, the map
$\SHR^r: \bitpoly{32} \to \bitpoly{32}$ extends $R$-linearly to
$R^{<32}[X]$ and commutes with multilinear extension.
\end{lemma}

This is an instance of Lemma~2.3 of the main paper: any $R$-linear
coordinate-wise map admits a virtual-column realization.
Concretely, a witness column $\bp{v}$ of bit-polynomials gives free
oracle access to $\SHR^r(\bp{v})$ via the same multilinear extension.
The implementation registers these as bit-op specs of type
$\mathsf{ShiftR}(r)$ on the underlying committed column.

Define
\begin{equation}\label{eq:rho-sigma}
\rho_{\sigma_0}(X) := X^{25} + X^{14}, \qquad
\rho_{\sigma_1}(X) := X^{15} + X^{13}, \qquad \text{in } \F_2[X].
\end{equation}

\begin{lemma}[{$\Q[X]$ lift of $\sigma_0/\sigma_1$}]\label{lem:small-sigma}
Let $W \in [0\mathbin{..}2^{32}-1]$ and $\bp{y} \in \bitpoly{32}$.
Then $\bp{y} = \widehat{\sigma_0(W)}$ if and only if there exists
$\bp{o} \in \bitpoly{32}$ such that
\begin{equation}\label{eq:lsig0-Q}
\bp{W} \cdot \rho_{\sigma_0}(X) \;+\; \SHR^{3}(\bp{W}) \;-\; \bp{y} \;-\; 2\,\bp{o}
\;\in\; (X^{32}-1) \qquad \text{in } \Q[X],
\end{equation}
and analogously for $\sigma_1$ via $\rho_{\sigma_1}$ and $\SHR^{10}$.
\end{lemma}

In the trace, $\SHR^{r}(\bp{W})$ is realized by a bit-op virtual column
over $\bp{W}$ (Lemma~\ref{lem:shr-virtual}), so it is not a committed
column. The slack $\bp{o}$ is a public bit-polynomial overflow witness
of bit-budget $1$ per coefficient.

\subsection{The functions $\Ch$ and $\Maj$ via Table~9 virtual residuals}
\label{sec:ch-maj}

Rather than committing the materialization columns
$B_1,B_2,B_3$ of the spec's Table~9 (each storing one $\Q[X]$ affine
combination meant to lie in $\bitpoly{32}$ per coefficient), we declare
those affine combinations as \emph{virtual binary-polynomial columns}
and pin them coefficient-wise via the booleanity sumcheck.

The $\Ch$ function is split into two AND-operand bit-polys
$\bp{u_{ef}} := e \AND f$ and $\bp{u_{\neg e,g}} := (\NOT e) \AND g$; since
$e$ and $\NOT e$ have disjoint bit-support, $\Ch(e,f,g) = u_{ef} + u_{\neg e,g}$
coefficient-wise in $\Z$.

An \emph{anchor} is the row index $k$ at which a row-local
constraint is centred: every cell access in the expression is written
as a forward shift $\bp{x}^{\downarrow\Delta}\rowidx{k} = \bp{x}\rowidx{k+\Delta}$
of $k$.\footnote{Forward-only because the protocol's row-shift virtual
columns are forward shifts; choosing the anchor at the smallest row
the constraint reads keeps every access on a real \texttt{ShiftSpec}.}
Anchoring at $k = t-2$ so all references are forward shifts, the three
virtual residuals are
\begin{align}
r_{\mathrm{ch}_1}[k]
  &:= e[k+2] + e[k+1] - 2\,u_{ef}[k+2], \label{eq:rch1}\\
r_{\mathrm{ch}_2}[k]
  &:= e[k+2] - e[k] + 2\,u_{\neg e,g}[k+2] + 2\,c_{\mathrm{ch}_2}[k],
        \label{eq:rch2}\\
r_{\mathrm{maj}}[k]
  &:= a[k] + a[k+1] + a[k+2] - 2\,\Maj[k+2] - 2\,c_{\mathrm{maj}}[k],
        \label{eq:rmaj}
\end{align}
where $c_{\mathrm{ch}_2}, c_{\mathrm{maj}} \in \bitpoly{32}$ are
\emph{tail compensators} discussed below. The booleanity sumcheck pins
each per-bit slice of these residuals into $\{0,1\}$, which forces every
coefficient of every residual into $\{0,1\}$ for free.

\paragraph{Why $r_{\mathrm{ch}_2}$ uses the alt-complement form.}
The natural Table~9 row 63 expression
$(\mathbf{1}_{32} - e[t]) + e[t-2] - 2\, u_{\neg e,g}[t]$ contains a
$\mathbf{1}_{32}$ constant; storing or accessing it virtually in the
sumcheck framework would require materializing a $32$-bit constant
column. The implementation observes that, per coefficient, the original
residue $r_{\mathrm{orig}}$ and the alt-complement residue
$r_{\mathrm{alt}}$ used in \eqref{eq:rch2} satisfy
$r_{\mathrm{orig}} + r_{\mathrm{alt}} = 1$, so they are bit-valid
simultaneously. Substituting the alt form lets us drop the
$\mathbf{1}_{32}$ constant from the virtual-source enum entirely.

\paragraph{Tail compensators.}
On every row $k$ with $k+2 < n$ the residuals collapse to clean
$\XOR$/$\Maj$-$\XOR$ patterns and lie in $\{0,1\}$ per coefficient. The
exceptions are the last two rows $k \in \{n-2, n-1\}$, where the
length-$2$ forward shifts pull off-trace zero-padding. Without
compensators:
\begin{itemize}
\item $r_{\mathrm{ch}_2}$ at $k \in \{n-2,n-1\}$ reduces to $-e[k]$,
which is in $\{0,-1\}$ — bit-invalid.
\item $r_{\mathrm{maj}}$ at $k = n-2$ reduces to $a[k]+a[k+1]$, which
is in $\{0,1,2\}$ — also bit-invalid.
\end{itemize}
The public compensators are therefore non-zero only on those two tail
rows, with values
\begin{equation}\label{eq:compensators}
c_{\mathrm{ch}_2}[k] := \begin{cases} e[k] & k+2 \ge n \\ 0 & \text{else}\end{cases}, \quad
c_{\mathrm{maj}}[k]  := \begin{cases} a[k] \AND a[k+1] & k+2 \ge n,\, k+1 < n \\ 0 & \text{else}.\end{cases}
\end{equation}
Plugging \eqref{eq:compensators} into \eqref{eq:rch2} and \eqref{eq:rmaj}
recovers $\{0,1\}$-valued residuals on every row. Both compensators are
deterministic functions of the public boundary cells $\bp{a}[n-2,n-1]$
and $\bp{e}[n-2,n-1]$ exposed by the digest pinning (see
\Cref{sec:cols}), so they are placed in the public input rather than
committed.

\subsection{Modular addition via the $(X-2)$ bridge}\label{sec:mod-add}

Recall from the main paper that, for a polynomial $f \in \Q[X]$,
membership $f \in (X-2)$ is equivalent to $f(2)=0$. We use this to
glue together $\bitpoly{32}$ values (whose integer image is $\bp{x}(2)$)
with $\Q$-valued accumulators.

\begin{lemma}[Modular addition lemma]\label{lem:mod-add}
Let $\bp{y}, \bp{z}_1,\dots,\bp{z}_M \in \bitpoly{32}$,
$\bp{x}_1,\dots,\bp{x}_L \in \bitpoly{32}$, and $\mu \in \Z_{\ge 0}$.
Suppose
\begin{equation}\label{eq:mod-add}
\bp{y} \;-\; \sum_{j=1}^{L} \bp{x}_j \;-\; \sum_{j=1}^{M} \bp{z}_j \;+\; 2 X^{31} \cdot \mu \;\in\; (X-2).
\end{equation}
Then
$\psi_2(\bp{y}) \equiv \sum_{j=1}^{L} \psi_2(\bp{x}_j) + \sum_{j=1}^{M} \psi_2(\bp{z}_j) \pmod{2^{32}}$.
The converse holds with $\mu \in [0\mathbin{..}L+M-1]$.
\end{lemma}

The substitution $2 X^{31}$ in place of $X^{32}$ lets us stay inside
the $\Q[X]$-ring of polynomials of degree~$<32$ used throughout the
implementation; the two coincide when evaluated at $X = 2$.

\begin{remark}[Bit-budget for $\mu$]\label{rem:mu-budget}
For each modular sum, the integer carry $\mu$ is bounded by the number
of summands. The implementation packs the five carries arising in the
SHA round (the message-schedule recurrence, the $a$- and $e$-updates,
and the two feed-forward additions) into disjoint bit-fields of a single
binary-polynomial column $\bp{W}_{\mu\text{-packed}}$, with the bit
layout
\begin{center}
\begin{tabular}{ll}
\toprule
\textbf{Bit range} & \textbf{Quantity} \\
\midrule
$0$--$1$  & $\mu_W$ \quad (msg-schedule carry, range $[0,3]$)\\
$2$--$4$  & $\mu_a$ \quad ($a$-update carry, range $[0,6]$)\\
$5$--$7$  & $\mu_e$ \quad ($e$-update carry, range $[0,5]$)\\
$8$       & $\mu_{ff,a}$ \quad (junction carry, range $\{0,1\}$)\\
$9$       & $\mu_{ff,e}$ \quad (junction carry, range $\{0,1\}$)\\
$10$--$31$ & must be $0$ (range-pinned by C\ref{con:mu-zero}). \\
\bottomrule
\end{tabular}
\end{center}
Each carry is read by a $\mathsf{ShiftR}$-virtual column on
$\bp{W}_{\mu\text{-packed}}$. Booleanity gives every bit of the packed
column the constraint $b \in \{0,1\}$ for free; for non-power-of-$2$
honest ranges (e.g.\ $\mu_a \le 6$, encoded in $3$ bits which would also
admit $\mu_a = 7$), violations are caught by the booleanity of the
$\bp{W}_{a}$ / $\bp{W}_{e}$ columns at the next row, since they would
force $a[k+4]$ or $e[k+4]$ to wrap negative.
\end{remark}
